\documentclass[12pt]{article}
\usepackage[utf8]{inputenc}
\usepackage[T1]{fontenc}
\usepackage{hyperref}
\title{Sarasira Language Resources for K-12 Curriculum Development}
\author{Compiled by Claude Sonnet 5.5 (claude-sonnet-5-5) for Andr\'as Kornai}
\date{}
\begin{document}
\maketitle

\section{Introduction}
Sarasira (ISO 639-3 zsa) is an Austronesian language of Morobe Province, Papua New Guinea, spoken in the upper Markham valley. Sources disagree on how many speakers it has and on how distinct it is from Adzera. Joshua Project, citing Ethnologue, gives about 1,200 speakers. Global Recordings Network files it as ``Adzera: Sarasira'' with a much larger population figure. Very little material appears to exist, and none of it is school material.

\section{Primary Educational Resources}
I found no primers, textbooks or school curricula for Sarasira in the search results.

\section{Linguistic and Reference Materials}
\begin{itemize}
\item Wikipedia, Adzera language. It lists Sarasira, Som, Pukpuk, Saseang and Sisuk as villages and notes that the classification of Sarasira relative to Adzera is unsettled. \url{https://en.wikipedia.org/wiki/Adzera_language}
\item PARADISEC catalogue. It holds field recordings of Markham-valley languages, including Sarasira speakers recorded in 1985 by Susanne Holzknecht. \url{https://catalog.paradisec.org.au/repository/SH2/ANUA520_5001}
\end{itemize}

\section{Religious and Community Resources}
\begin{itemize}
\item Joshua Project, Sirasira in Papua New Guinea. It lists the Bible translation status as ``Translation Started''. \url{https://joshuaproject.net/people_groups/14947}
\end{itemize}

\section{Digital and Audio Resources}
\begin{itemize}
\item Global Recordings Network, Sarasira language page. \url{https://globalrecordings.net/en/language/6697}
\end{itemize}

\section{Curriculum Development Context}
No evidence was found that Sarasira is taught in schools. A curriculum would have to be built from scratch with community speakers. The PARADISEC recordings and any wordlists related to Adzera are the only documentary starting points found. Any Adzera material should be checked carefully before it is treated as Sarasira.

\end{document}